\section{Discussion}

\subsection{Scheme Selection}
NGP minimizes deposit-kernel time but exhibits the highest spectral noise (0.61) and largest 500-step energy drift.
CIC provides the best accuracy--performance balance on CPU: fastest deposit-kernel time at moderate $N_p$, noise reduced by 36\%, and energy drift an order of magnitude lower than NGP.
TSC further reduces noise to 0.26 at roughly $1.5\times$ the CIC kernel time.
Esirkepov maintains trajectory-based charge conservation in full PIC cycles but remains $\sim 7\times$ slower than CIC at $N_p=10^5$ due to per-particle weight computation along crossed cells.

\subsection{Memory Locality and Data Layout}
Sorting improves spatial locality by clustering particles that share grid neighborhoods, but its $\mathcal{O}(N_p \log N_p)$ cost exceeds the deposit kernel at $N_p \leq 10^5$ when timed per deposition call.
In full PIC cycles, sort cost can be amortized if particle order is preserved across steps or triggered periodically rather than every step.
SoA and AoS achieve comparable deposit-kernel throughput on Apple Silicon; SoA's advantage in unit-stride attribute reads is partially offset by the higher cost of SoA reordering.
The choice of layout should therefore be guided by the full timestep profile, not the deposit kernel alone.

\subsection{Parallelism Limits}
OpenMP privatization removes atomic contention and yields a $4.4\times$ speedup from one to eight threads for CIC, but scaling flattens beyond two threads.
Each thread allocates a full-grid private buffer ($T \times N_x N_y$ doubles) and merges through a critical section, so merge cost and memory bandwidth become the bottleneck.
GPU atomics versus privatized tiles remain to be benchmarked on NVIDIA hardware where higher memory bandwidth and many lightweight threads may shift this balance.

\section{Conclusion}
We presented a unified, reproducible framework for evaluating charge deposition in 2D electrostatic PIC simulations.
The codebase implements four deposition schemes behind a common runtime interface, separates sort from deposit timing, and couples performance benchmarks with charge, energy, spectral-noise, and Langmuir-wave validation.
On Apple Silicon, CIC offers the strongest accuracy--performance trade-off for CPU production runs; layout choice depends on whether sort overhead or attribute-access patterns dominate the full timestep.
Esirkepov is preferable when trajectory-based charge conservation is paramount and particle counts are moderate.
NGP is viable only when raw kernel speed outweighs noise tolerance.

\section{Future Work}
Planned extensions include:
\begin{itemize}
    \item CUDA benchmarking of atomic and privatized-tile kernels on Linux clusters with cuFFT-coupled Poisson solves.
    \item MPI domain decomposition with ghost-cell exchange to scale beyond shared memory.
    \item 3D extension with face-aligned sorting and wider stencils.
    \item Electromagnetic PIC with charge-conserving current deposition~\cite{decyk1996,wilke2017}.
    \item Incremental and GPU-resident sorting to reduce per-step reorder overhead.
    \item Absolute Langmuir-frequency calibration with corrected 2D plasma dispersion.
\end{itemize}
